\section{Finite interleaving}

\import{set/finite-sequence.tex}

This module interleaves a sequence of entries with a sequence of values to be
inserted between them.  The indexing results are stated separately so that
later finite-chain arguments can use the exact even and odd positions without
introducing division or parity arithmetic.

\begin{lemma}\label{natural_double_injective}
    Let $i,j$ be natural numbers.
    Suppose $i+i=j+j$.
    Then $i=j$.
\end{lemma}
\begin{proof}
    Let $C=\{k\in\naturals\mid
        \text{for every $l\in\naturals$ such that $k+k=l+l$
        we have $k=l$}\}$.
    Show that $\emptyset\in C$.
    \begin{subproof}
        $\emptyset$ is a natural number by \obv{omega_empty}.
        Show that for every $l\in\naturals$ such that
            $\emptyset+\emptyset=l+l$ we have $\emptyset=l$.
        \begin{subproof}
            Fix $l$.
            Assume $l\in\naturals$ and $\emptyset+\emptyset=l+l$.
            Suppose not.
            Then $l\neq\emptyset$.
            Take a natural number $q$ such that $l=\suc{q}$
                by \obv{natural_zero_or_successor}.
            $l+l=\suc{\suc{q+q}}$
                by \auto{num_addition_successor_left,num_addition_successor_right}.
            $\emptyset+\emptyset=\emptyset$
                by \obv{num_addition_zero_right}.
            Contradiction by \obv{suc_neq_emptyset}.
        \end{subproof}
    \end{subproof}
    Show that for every $k\in C$ we have $\suc{k}\in C$.
    \begin{subproof}
        Fix $k\in C$.
        Hence $\suc{k}$ is a natural number by \obv{omega_successor}.
        Show that for every $l\in\naturals$ such that
            $\suc{k}+\suc{k}=l+l$ we have $\suc{k}=l$.
        \begin{subproof}
            Fix $l$.
            Assume $l\in\naturals$ and $\suc{k}+\suc{k}=l+l$.
            $l\neq\emptyset$
                by \auto{num_addition_zero_right,num_addition_successor_left,num_addition_successor_right,suc_neq_emptyset}.
            Take a natural number $q$ such that $l=\suc{q}$
                by \obv{natural_zero_or_successor}.
            $\suc{k}+\suc{k}=\suc{\suc{k+k}}$ and
                $l+l=\suc{\suc{q+q}}$
                by \auto{num_addition_successor_left,num_addition_successor_right}.
            Hence $k+k=q+q$ by \obv{suc_injective}.
            Hence $k=q$.
        \end{subproof}
    \end{subproof}
    Hence $C$ is an inductive set by \obv{inductive_set}.
    Hence every natural number is an element of $C$
        by \obv{omega_minimal}.
    Thus $i\in C$.
\end{proof}

\begin{lemma}\label{natural_double_not_successor_double}
    Let $i,j$ be natural numbers.
    Then $i+i\neq\suc{j+j}$.
\end{lemma}
\begin{proof}
    Let $C=\{k\in\naturals\mid
        \text{for every $l\in\naturals$ we have
        $k+k\neq\suc{l+l}$}\}$.
    Show that $\emptyset\in C$.
    \begin{subproof}
        $\emptyset$ is a natural number by \obv{omega_empty}.
        For every $l\in\naturals$ we have
            $\emptyset+\emptyset\neq\suc{l+l}$
            by \auto{num_addition_zero_right,suc_neq_emptyset}.
    \end{subproof}
    Show that for every $k\in C$ we have $\suc{k}\in C$.
    \begin{subproof}
        Fix $k\in C$.
        Hence $\suc{k}$ is a natural number by \obv{omega_successor}.
        Show that for every $l\in\naturals$ we have
            $\suc{k}+\suc{k}\neq\suc{l+l}$.
        \begin{subproof}
            Fix $l\in\naturals$.
            Suppose $\suc{k}+\suc{k}=\suc{l+l}$.
            $l\neq\emptyset$
                by \auto{num_addition_zero_right,num_addition_successor_left,num_addition_successor_right,suc_injective,suc_neq_emptyset}.
            Take a natural number $q$ such that $l=\suc{q}$
                by \obv{natural_zero_or_successor}.
            $\suc{k}+\suc{k}=\suc{\suc{k+k}}$ and
                $\suc{l+l}=\suc{\suc{\suc{q+q}}}$
                by \auto{num_addition_successor_left,num_addition_successor_right}.
            Hence $k+k=\suc{q+q}$ by \obv{suc_injective}.
            $k+k\neq\suc{q+q}$.
            Contradiction.
        \end{subproof}
    \end{subproof}
    Hence $C$ is an inductive set by \obv{inductive_set}.
    Hence every natural number is an element of $C$
        by \obv{omega_minimal}.
    Thus $i\in C$.
\end{proof}

\begin{definition}\label{finite_even_indices}
    $\finiteEvenIndices{r}=\{i+i\mid i\in\suc{r}\}$.
\end{definition}

\begin{definition}\label{finite_odd_indices}
    $\finiteOddIndices{r}=\{\suc{i+i}\mid i\in r\}$.
\end{definition}

\begin{theorem}\label{finite_double_index_partition}
    Let $r$ be a natural number.
    Then $\suc{r+r}=\finiteEvenIndices{r}\union\finiteOddIndices{r}$.
\end{theorem}
\begin{proof}
    Let $C=\{n\in\naturals\mid
        \suc{n+n}=\finiteEvenIndices{n}\union\finiteOddIndices{n}\}$.
    Show that $\emptyset\in C$.
    \begin{subproof}
        $\emptyset$ is a natural number by \obv{omega_empty}.
        $\emptyset+\emptyset=\emptyset$
            by \obv{num_addition_zero_right}.
        $\suc{\emptyset+\emptyset}=\{\emptyset+\emptyset\}$
            by \obv{suc,singleton_iff,cons_iff,setext}.
        $\emptyset\in\suc{\emptyset}$ by \obv{suc_intro_self}.
        $\finiteEvenIndices{\emptyset}
            =\{\emptyset+\emptyset\}$
            by \auto{finite_even_indices,suc_iff,emptyset,setext}.
        $\finiteOddIndices{\emptyset}=\emptyset$
            by \auto{finite_odd_indices,emptyset,setext}.
        Follows by \auto{union_emptyset}.
    \end{subproof}
    Show that for every $n\in C$ we have $\suc{n}\in C$.
    \begin{subproof}
        Fix $n\in C$.
        Hence $\suc{n}$ is a natural number by \obv{omega_successor}.
        $\suc{n}+\suc{n}=\suc{\suc{n+n}}$
            by \auto{num_addition_successor_left,num_addition_successor_right}.
        Show that
            $\suc{\suc{n}+\suc{n}}
                =\finiteEvenIndices{\suc{n}}
                    \union\finiteOddIndices{\suc{n}}$.
        \begin{subproof}
            Show that every element of
                $\suc{\suc{n}+\suc{n}}$ is an element of
                $\finiteEvenIndices{\suc{n}}
                    \union\finiteOddIndices{\suc{n}}$.
            \begin{subproof}
                Fix $q\in\suc{\suc{n}+\suc{n}}$.
                We have $q\in\suc{n+n}$ or $q=\suc{n+n}$ or
                    $q=\suc{\suc{n+n}}$
                    by \auto{suc_iff}.
                $\suc{n+n}
                    =\finiteEvenIndices{n}\union\finiteOddIndices{n}$.
                $\finiteEvenIndices{n}\subseteq
                    \finiteEvenIndices{\suc{n}}$ and
                    $\finiteOddIndices{n}\subseteq
                    \finiteOddIndices{\suc{n}}$
                    by \auto{finite_even_indices,finite_odd_indices,suc_iff,subseteq}.
                $\suc{n+n}\in\finiteOddIndices{\suc{n}}$
                    by \obv{finite_odd_indices,suc_intro_self}.
                $\suc{\suc{n+n}}=\suc{n}+\suc{n}$ and
                    $\suc{n}\in\suc{\suc{n}}$ by \obv{suc_intro_self}.
                Hence $\suc{\suc{n+n}}
                    \in\finiteEvenIndices{\suc{n}}$
                    by \auto{finite_even_indices}.
                Follows by \auto{union_iff,subseteq}.
            \end{subproof}
            Show that every element of
                $\finiteEvenIndices{\suc{n}}
                    \union\finiteOddIndices{\suc{n}}$ is an element of
                $\suc{\suc{n}+\suc{n}}$.
            \begin{subproof}
                Follows by
                    \auto{finite_even_indices,finite_odd_indices,union_iff,suc_iff,ordinal_elem_implies_subseteq,natural_number_is_ordinal,ordinal_subseteq_implies_elem_or_eq,num_addition_on_naturals,num_addition_successor_left,num_addition_successor_right}.
            \end{subproof}
            Follows by set extensionality.
        \end{subproof}
    \end{subproof}
    Hence $C$ is an inductive set by \obv{inductive_set}.
    Hence every natural number is an element of $C$
        by \obv{omega_minimal}.
    Thus $r\in C$.
\end{proof}

\begin{definition}\label{finite_interleave}
    $\finiteInterleave{c}{m}=
        \{(i+i,c(i))\mid i\in\dom{c}\}\union
        \{(\suc{i+i},m(i))\mid i\in\dom{m}\}$.
\end{definition}

\begin{proposition}\label{finite_interleave_function}
    Let $r$ be a natural number.
    Suppose $\dom{c}=\suc{r}$ and $\dom{m}=r$.
    Then $\finiteInterleave{c}{m}\in
        \funs{\suc{r+r}}{\ran{\finiteInterleave{c}{m}}}$.
\end{proposition}
\begin{proof}
    Let $d=\finiteInterleave{c}{m}$.
    Let $E=\{(i+i,c(i))\mid i\in\dom{c}\}$.
    Let $O=\{(\suc{i+i},m(i))\mid i\in\dom{m}\}$.
    Then $d=E\union O$ by \obv{finite_interleave}.
    $d$ is a relation
        by \auto{finite_interleave,relation,union_iff,pair_eq_iff}.
    Show that $d$ is right-unique.
    \begin{subproof}
        Show for all $j,x,y$ such that $(j,x),(j,y)\in d$ we have $x=y$.
        \begin{subproof}
            Fix $j,x,y$.
            Assume $(j,x),(j,y)\in d$.
            $(j,x)\in E$ or $(j,x)\in O$ by \auto{union_iff}.
            $(j,y)\in E$ or $(j,y)\in O$ by \auto{union_iff}.
            \begin{byCase}
            \caseOf{$(j,x)\in E$.}
                Take $i\in\dom{c}$ such that
                    $(j,x)=(i+i,c(i))$.
                Hence $i\in\suc{r}$.
                Thus $i$ is a natural number
                    by \auto{omega_successor,naturals_is_transitive_set,transitiveset_iff_subseteq,subseteq}.
                \begin{byCase}
                \caseOf{$(j,y)\in E$.}
                    Take $k\in\dom{c}$ such that
                        $(j,y)=(k+k,c(k))$.
                    Hence $k\in\suc{r}$.
                    Thus $k$ is a natural number
                        by \auto{omega_successor,naturals_is_transitive_set,transitiveset_iff_subseteq,subseteq}.
                    $i+i=k+k$ by \auto{pair_eq_iff}.
                    Hence $i=k$ by \obv{natural_double_injective}.
                    Follows by \auto{pair_eq_iff}.
                \caseOf{$(j,y)\notin E$.}
                    Then $(j,y)\in O$.
                    Take $k\in\dom{m}$ such that
                        $(j,y)=(\suc{k+k},m(k))$.
                    Hence $k\in r$.
                    Thus $k$ is a natural number
                        by \obv{naturals_is_transitive_set,transitiveset_iff_subseteq,subseteq}.
                    $i+i=\suc{k+k}$ by \auto{pair_eq_iff}.
                    Contradiction by \obv{natural_double_not_successor_double}.
                \end{byCase}
            \caseOf{$(j,x)\notin E$.}
                Then $(j,x)\in O$.
                Take $i\in\dom{m}$ such that
                    $(j,x)=(\suc{i+i},m(i))$.
                Hence $i\in r$.
                Thus $i$ is a natural number
                    by \obv{naturals_is_transitive_set,transitiveset_iff_subseteq,subseteq}.
                \begin{byCase}
                \caseOf{$(j,y)\in E$.}
                    Take $k\in\dom{c}$ such that
                        $(j,y)=(k+k,c(k))$.
                    Hence $k\in\suc{r}$.
                    Thus $k$ is a natural number
                        by \auto{omega_successor,naturals_is_transitive_set,transitiveset_iff_subseteq,subseteq}.
                    $k+k=\suc{i+i}$ by \auto{pair_eq_iff}.
                    Contradiction by \obv{natural_double_not_successor_double}.
                \caseOf{$(j,y)\notin E$.}
                    Then $(j,y)\in O$.
                    Take $k\in\dom{m}$ such that
                        $(j,y)=(\suc{k+k},m(k))$.
                    Hence $k\in r$.
                    Thus $k$ is a natural number
                        by \obv{naturals_is_transitive_set,transitiveset_iff_subseteq,subseteq}.
                    $i+i=k+k$ by \auto{pair_eq_iff,suc_injective}.
                    Hence $i=k$ by \obv{natural_double_injective}.
                    Follows by \auto{pair_eq_iff}.
                \end{byCase}
            \end{byCase}
        \end{subproof}
        Follows by \obv{rightunique}.
    \end{subproof}
    Show that $\dom{d}=\suc{r+r}$.
    \begin{subproof}
        Show that $\dom{d}\subseteq\suc{r+r}$.
        \begin{subproof}
            It suffices to show that every element of $\dom{d}$ is an element
                of $\suc{r+r}$ by \obv{subseteq}.
            Fix $j\in\dom{d}$.
            Take $x$ such that $(j,x)\in d$ by \obv{dom_iff}.
            $(j,x)\in E$ or $(j,x)\in O$ by \auto{union_iff}.
            Hence $j\in\finiteEvenIndices{r}$ or
                $j\in\finiteOddIndices{r}$
                by \auto{finite_even_indices,finite_odd_indices,pair_eq_iff}.
            Thus $j\in\finiteEvenIndices{r}\union\finiteOddIndices{r}$
                by \obv{union_iff}.
            Follows by \obv{finite_double_index_partition}.
        \end{subproof}
        Show that $\suc{r+r}\subseteq\dom{d}$.
        \begin{subproof}
            It suffices to show that every element of $\suc{r+r}$ is an
                element of $\dom{d}$ by \obv{subseteq}.
            Fix $j\in\suc{r+r}$.
            $j\in\finiteEvenIndices{r}\union\finiteOddIndices{r}$
                by \obv{finite_double_index_partition}.
            Follows by
                \auto{finite_even_indices,finite_odd_indices,finite_interleave,union_iff,dom_iff}.
        \end{subproof}
        Follows by \obv{subseteq_antisymmetric}.
    \end{subproof}
    Hence $d\in\funs{\suc{r+r}}{\ran{d}}$
        by \auto{function_exact_type_to_range}.
\end{proof}

\begin{proposition}\label{finite_interleave_domain}
    Let $r$ be a natural number.
    Suppose $\dom{c}=\suc{r}$ and $\dom{m}=r$.
    Then $\dom{\finiteInterleave{c}{m}}=\suc{r+r}$.
\end{proposition}
\begin{proof}
    $\finiteInterleave{c}{m}\in
        \funs{\suc{r+r}}{\ran{\finiteInterleave{c}{m}}}$
        by \obv{finite_interleave_function}.
    Follows by \obv{funs}.
\end{proof}

\begin{proposition}\label{finite_interleave_even}
    Let $r$ be a natural number.
    Suppose $\dom{c}=\suc{r}$ and $\dom{m}=r$.
    Let $i\in\suc{r}$.
    Then $\apply{\finiteInterleave{c}{m}}{i+i}=c(i)$.
\end{proposition}
\begin{proof}
    Hence $i\in\dom{c}$.
    $(i+i,c(i))\in\{(j+j,c(j))\mid j\in\dom{c}\}$
        by \auto.
    Hence $(i+i,c(i))\in\finiteInterleave{c}{m}$
        by \auto{finite_interleave,union_iff}.
    $\finiteInterleave{c}{m}$ is a function
        by \auto{finite_interleave_function,funs_is_function}.
    Follows by \obv{function_apply_intro}.
\end{proof}

\begin{proposition}\label{finite_interleave_odd}
    Let $r$ be a natural number.
    Suppose $\dom{c}=\suc{r}$ and $\dom{m}=r$.
    Let $i\in r$.
    Then $\apply{\finiteInterleave{c}{m}}{\suc{i+i}}=m(i)$.
\end{proposition}
\begin{proof}
    Hence $i\in\dom{m}$.
    $(\suc{i+i},m(i))
        \in\{(\suc{j+j},m(j))\mid j\in\dom{m}\}$
        by \auto.
    Hence $(\suc{i+i},m(i))\in\finiteInterleave{c}{m}$
        by \auto{finite_interleave,union_iff}.
    $\finiteInterleave{c}{m}$ is a function
        by \auto{finite_interleave_function,funs_is_function}.
    Follows by \obv{function_apply_intro}.
\end{proof}

\begin{proposition}\label{finite_interleave_typed}
    Let $r$ be a natural number.
    Suppose $c\in\funs{\suc{r}}{A}$ and $m\in\funs{r}{A}$.
    Then $\finiteInterleave{c}{m}\in\funs{\suc{r+r}}{A}$.
\end{proposition}
\begin{proof}
    $\dom{c}=\suc{r}$ and $\dom{m}=r$ and
        $c$ and $m$ are functions by \obv{funs_elim,funs_is_function}.
    Let $d=\finiteInterleave{c}{m}$.
    $d\in\funs{\suc{r+r}}{\ran{d}}$
        by \auto{finite_interleave_function}.
    $\ran{c}\subseteq A$ and $\ran{m}\subseteq A$ by \obv{funs_ran}.
    Show that for every $y\in\ran{d}$ we have
        $y\in\ran{c}$ or $y\in\ran{m}$.
    \begin{subproof}
        Fix $y\in\ran{d}$.
        Take $j$ such that $(j,y)\in d$ by \obv{ran_iff}.
        $(j,y)\in\{(i+i,c(i))\mid i\in\dom{c}\}$ or
            $(j,y)\in\{(\suc{i+i},m(i))\mid i\in\dom{m}\}$
            by \auto{finite_interleave,union_iff}.
        \begin{byCase}
        \caseOf{$(j,y)\in\{(i+i,c(i))\mid i\in\dom{c}\}$.}
            Take $i\in\dom{c}$ such that $(j,y)=(i+i,c(i))$
                by \auto.
            Hence $y=c(i)$ by \auto{pair_eq_iff}.
            Thus $y\in\ran{c}$ by \auto{fun_ran_iff}.
        \caseOf{$(j,y)\notin\{(i+i,c(i))\mid i\in\dom{c}\}$.}
            Then $(j,y)\in
                \{(\suc{i+i},m(i))\mid i\in\dom{m}\}$.
            Take $i\in\dom{m}$ such that
                $(j,y)=(\suc{i+i},m(i))$ by \auto.
            Hence $y=m(i)$ by \auto{pair_eq_iff}.
            Thus $y\in\ran{m}$ by \auto{fun_ran_iff}.
        \end{byCase}
    \end{subproof}
    Hence $\ran{d}\subseteq A$ by \obv{subseteq}.
    Hence $d\in\funs{\suc{r+r}}{A}$ by \obv{funs_weaken_codom}.
\end{proof}
